% Laget automatisk av elevark/lag_elevark.py fra sim/maaleusikkerhet.html. Endringer her blir overskrevet.
\documentclass[11pt]{article}
\usepackage{elevark}
\begin{document}
\elevarkhode{Fysikk 1 · Introduksjon til fysikk}{Måleusikkerhet}{https://fysikklab.no/sim/maaleusikkerhet.html}{Simuleringen}

Du skal måle svingetiden til en pendel med stoppeklokke. Hver gang du trykker, kommer du litt for tidlig eller litt for sent. Hvor nøyaktig kan du bestemme svingetiden, og hvordan kan du gjøre målingen bedre?

\section*{Oppgaver}

\begin{oppgaveliste}
\oppgave \fase{test}{Mål}Mål svingetiden én periode om gangen, ti ganger. Hva er middelverdien? Hvor stort er det største avviket fra middelverdien? Skriv resultatet på formen $T = (\text{middelverdi} \pm \text{usikkerhet})\,\mathrm{s}$ med fornuftig antall desimaler.
\svar
\oppgave \fase{gjett}{Gjett}Hva tror du skjer med spredningen hvis du i stedet måler tiden for 10 perioder og deler på 10? Test.
\svar
\oppgave \fase{forklar}{Forklar}Hvorfor blir spredningen omtrent ti ganger mindre når du måler ti perioder om gangen? Feilen i hvert trykk er jo like stor.
\svar
\oppgave \fase{test}{Undersøk}Mål mange ganger, for eksempel 10, 40 og 100 målinger med én periode. Hva skjer med standardavviket? Hva skjer med usikkerheten i middelverdien?
\svar
\oppgave \fase{test}{Undersøk}Slå på den systematiske feilen og mål mange ganger. Slå så på «Vis sann verdi». Hjelper det å måle mange ganger? Hva kan du gjøre for å redusere virkningen av feilen?
\svar
\oppgave \fase{forklar}{Forklar}Forklar forskjellen på tilfeldige og systematiske feil. Gi et eksempel på hver fra et forsøk du har gjort selv.
\svar
\oppgave \fase{utfordring}{Utfordring}Svingetiden er $T = 2π\sqrt{L/g}$. Bruk middelverdien din til å regne ut $g$. Hvor stor blir den relative usikkerheten i $g$ når den relative usikkerheten i $T$ er 1~\%?
\svar
\end{oppgaveliste}

\end{document}
